(1) 取剪应变 \(\gamma=\partial_xu\)，本构变为 \((1-i\omega\tau)\sigma=-i\omega\eta\gamma\)，故复剪切模量 \(G^*(\omega)=\sigma/\gamma=(-i\omega\eta)/(1-i\omega\tau)\)。采用 \(e^{-i\omega t}\) 约定，其损耗项符号与正耗散一致。
(2) 动量方程给出 \(-\rho\omega^2u=ik\sigma\)，联立 \(\gamma=iku\) 得 \(k^2=\rho(\tau\omega^2+i\omega)/\eta\)。选择 \(\Re k>0,\Im k>0\) 的根使沿 \(+x\) 传播波随距离衰减，分支选择不能只看平方根的绝对值。
(3) 低频 \(k\simeq(1+i)\sqrt{\rho\omega/(2\eta)}\)，为黏性剪切扩散；高频 \(k\simeq\omega\sqrt{\rho\tau/\eta}+i\sqrt{\rho/(\eta\tau)}/2\)，相速度趋 \(c_\infty=\sqrt{\eta/(\rho\tau)}\)，对应高频剪切模量 \(G_\infty=\eta/\tau\)。
(4) 本构乘 \(\sigma/\eta\) 可写 \(\sigma\dot\gamma=(\sigma^2/\eta)+(\tau/\eta)\sigma\dot\sigma\)，即输入功率密度分解为不可逆耗散 \(\sigma^2/\eta\ge0\) 与弹性储能变化 \(d(\tau\sigma^2/(2\eta))/dt\)。Maxwell材料长期流动，区别于具有静态剪切模量的Kelvin--Voigt固体。
